% ==============================================================================
% Section 4: Experiments - Thermogeometric Lift Test
% ==============================================================================
% Draft for CIFAR-10 validation of the Unified Scaling Law
% Author: Leo Liu / ThermoRG-NN Team
% Date: 2026-03-28
% Status: DRAFT - awaiting training results
% ==============================================================================

\section{Experiments: Thermogeometric Lift Test}
\label{sec:experiments}

\subsection{Dataset and Setup}
\label{sec:dataset_setup}

We validate the Unified Scaling Law on CIFAR-10 \citep{cifar10}, chosen for its 
tractable experimental timescale and well-characterized manifold structure. 
CIFAR-10 consists of 60,000 $32\times32$ RGB images across 10 balanced classes, 
providing sufficient complexity to reveal thermogeometric effects while remaining 
computationally feasible for comprehensive architecture sweeps.

Training proceeds with standard data augmentation: random horizontal flips, 
random cropping with 4-pixel padding, and normalization. We use SGD with momentum 
0.9 and weight decay $5\times10^{-4}$. The initial learning rate is 0.1, 
reduced by a factor of 10 at epochs 60 and 90 (for 150-epoch training) or 
at epoch 20 (for 30-epoch screening).

\subsection{Architecture Control Groups}
\label{sec:control_groups}

To isolate thermogeometric effects from confounding factors (depth, width, 
parameter count), we design four control groups spanning 15 architectures:

\begin{itemize}
  \item \textbf{Group 1 (Thermogeometric Optimal):} 
    Architectures satisfying both constraints:
    \begin{align}
      \mathcal{J}_{\text{topo}} &\leq \epsilon_{\text{topo}} \quad \text{(C1)} \\
      \tilde{T}_{\text{eff}} &\leq \xi_{\text{opt}} \, T_c \quad \text{(C2)}
    \end{align}
    This group includes ResNet-18-TGR, ResNet-34-TGR, WideResNet-28-10-TGR, 
    and DenseNet-40-TGR, where TGR denotes the Thermogeometric Reconfiguration 
    layer placement and normalization strategy.
    
  \item \textbf{Group 2 (Topology Destroyer):}
    Architectures with identical parameter counts to Group 1 but with 
    $\mathcal{J}_{\text{topo}} > \epsilon_{\text{topo}}$ via narrow bottleneck 
    channels (reduction factor 0.25 instead of 0.5). This breaks C1 while 
    preserving FLOP count. Examples: ResNet-18-Narrow, WideResNet-28-10-Narrow.
    
  \item \textbf{Group 3 (Thermal Boiling Furnace):}
    Architectures with intact C1 ($\mathcal{J}_{\text{topo}}$ within tolerance) 
    but broken C2 via raw ReLU activations without batch normalization or 
    Thermogeometric Reconfiguration. The effective temperature $\tilde{T}_{\text{eff}}$ 
    exceeds $\xi_{\text{opt}} T_c$, simulating thermal disorder. 
    Examples: ResNet-18-Furnace, DenseNet-40-Furnace.
    
  \item \textbf{Group 4 (Traditional Baselines):}
    Standard architectures without thermogeometric modifications, serving as 
    industry anchors: ResNet-18 \citep{resnet}, VGG-11 \citep{vgg}, and DenseNet-40 
    \citep{densenet}.
\end{itemize}

This $2\times2$ design (C1 $\times$ C2) enables causal isolation of each 
constraint's contribution to grokking dynamics and final performance.

\subsection{Two-Phase Evaluation Protocol}
\label{sec:evaluation_protocol}

We employ a sequential screening protocol to balance statistical power with 
computational cost:

\begin{itemize}
  \item \textbf{Phase A (30 epochs, all 15 architectures):} 
    Rapid screening for early grokking dynamics. We track test loss curves 
    to identify architectures exhibiting phase transitions in the 
    $[15, 30]$ epoch window. This phase filters architectures for detailed 
    analysis while maintaining Bonferroni-corrected significance thresholds.
    
  \item \textbf{Phase B (150 epochs, 6 selected architectures):}
    Full training of top candidates from Phase A (2 per group, selected by 
    lowest Phase A test loss). This phase validates scaling law correlations 
    and provides grokking transition timing measurements with $\pm 1$ epoch 
    resolution.
\end{itemize}

Each configuration is trained with 3 random seeds (std $\pm 0.3\%$ on final 
accuracy), yielding $N=45$ runs in Phase A and $N=18$ runs in Phase B.

\subsection{Metrics and Analysis}
\label{sec:metrics}

We evaluate the Unified Scaling Law through three complementary metrics:

\begin{enumerate}
  \item \textbf{Spearman Rank Correlation ($\rho_s$):}
    Between predicted complexity exponent $\alpha$ and actual test accuracy 
    after 150 epochs. We hypothesize $\rho_s > 0.7$ for Group 1 architectures 
    satisfying both constraints, and $\rho_s < 0.3$ for Group 2/3 where 
    constraints are violated.
    
  \item \textbf{Grokking Phase Transition Timing ($\tau_g$):}
    The epoch at which $\partial \mathcal{L}_{\text{test}}/\partial t$ reaches 
    maximum magnitude, marking the onset of phase transition from memorization 
    to generalization. We predict $\tau_g \propto \exp(\Delta S/k_B)$ where 
    $\Delta S$ is the entropy gap between initial and final states.
    
  \item \textbf{Pareto-Optimal Frontier in $(T_c, \gamma_c)$ Space:}
    The set of architectures not dominated by any other in both landscape 
    flatness ($T_c$) and kinetic mobility ($\gamma_c$). We expect Group 1 
    architectures to cluster near the Pareto frontier, while Group 2/3 
    architectures scatter into the infeasible region.
\end{enumerate}

\subsection{Preliminary Results}
\label{sec:preliminary_results}

\textit{[This section will be populated with experimental results. Expected 
figures: (a) Spearman correlation scatter with $\rho_s = 0.XX$ ($p < 0.XX$), 
(b) grokking dynamics comparing G1 vs G3 with shaded standard error regions, 
(c) Pareto front in $(T_c, \gamma_c)$ space with feasible region shaded.]}

\subsection{Discussion}
\label{sec:discussion}

The thermogeometric perspective predicts that architectures satisfying both 
C1 (topological efficiency) and C2 (thermal equilibrium) will exhibit:
\begin{itemize}
  \item Earlier grokking transitions due to lower effective entropy barriers
  \item Higher final accuracy due to improved landscape traversability  
  \item Tighter correlation between predicted $\alpha$ and realized performance
\end{itemize}

Conversely, architectures in Group 2 (broken C1) should suffer from 
frustrated gradient flow, while Group 3 (broken C2) should exhibit delayed 
grokking due to kinetic trapping in sharp minima.

% ==============================================================================
% References (add to main bibliography)
% ==============================================================================
% \citep{cifar10}    - CIFAR-10 dataset
% \citep{resnet}     - He et al., Deep Residual Learning
% \citep{vgg}        - Simonyan & Zisserman, VGG Networks  
% \citep{densenet}   - Huang et al., DenseNet
% ==============================================================================
